\documentclass[a4paper,12pt]{article}

% --- 俄语支持与必备宏包 ---
\usepackage[T2A]{fontenc}
\usepackage[utf8]{inputenc}
\usepackage[russian]{babel}
\usepackage{amsmath, amssymb, amsthm}
\usepackage{geometry}
\usepackage{booktabs}
\usepackage{float}
\usepackage{listings} % 用于插入代码
\usepackage{xcolor}

% 代码块样式设置
\lstset{
    language=Python,
    basicstyle=\ttfamily\small,
    keywordstyle=\color{blue},
    stringstyle=\color{red},
    commentstyle=\color{green!60!black},
    numbers=left,
    numberstyle=\tiny\color{gray},
    stepnumber=1,
    frame=single,
    breaklines=true,
    captionpos=b
}

\geometry{left=20mm, right=20mm, top=20mm, bottom=20mm}

\begin{document}

\begin{titlepage}
    \centering
    \textbf{Университет ИТМО} \\
    \vspace{1em}
    Факультет Программной Инженерии и Компьютерной Техники\\
    \vspace{1em}
    \vspace{4cm}
    \textbf{\Large Расчётно-графическая работа №2} \\
    \textbf{Дисциплина: Математическая статистика} \\
    \vspace{4cm}
    \textbf{Тема: Проверка статистических гипотез} \\
    \vfill
    Выполнил: Чжун Ц.
              Су Л.\\
    Вариант: А-4 \\
    2026
\end{titlepage}

\section{Программный код (Python)}
Для проведения всех статистических расчетов был использован язык программирования Python и библиотека \texttt{scipy.stats}. Код учитывает необходимость объединения интервалов для критерия Пирсона.

\begin{lstlisting}[caption=Реализация статистических критериев]

\end{lstlisting}

\section{Результаты вычислений}

\subsection{Сравнение средних $X_1$ и $X_2$ (Этап 4.2 и 4.4)}
\begin{itemize}
    \item \textbf{Параметрический $t$-критерий:} $h_{набл} = [0.3593]$, $p\text{-value} = [0.7197]$.
    \item \textbf{Критерий Манна-Уитни:} $U_{набл} = [5785.0000]$, $p\text{-value} = [0.709275]$.
    \item \textbf{Вывод:} [Напишите: "Средние равны" или "Средние различаются"] на уровне значимости $\alpha=0.05$. Выводы обоих критериев [совпали/не совпали].
\end{itemize}

\subsection{Параметр нормального распределения $X_3$ (Этап 4.3)}
Проверялась гипотеза $H_0: \mu = 95.0$.
\begin{itemize}
    \item Наблюдаемое значение: $h = [-1.4767]$.
    \item $p\text{-value} = [0.1427]$.
    \item \textbf{Результат:} Гипотеза $H_0$ [отвергается / не отвергается].
\end{itemize}

\subsection{Критерий согласия Пирсона для $X_4$ (Этап 4.5)}
Проверялось соответствие экспоненциальному закону $Exp(0.1)$.
\begin{itemize}
    \item Число степеней свободы $df = [5]$.
    \item Статистика $\chi^2_{набл} = [13.2211]$, $p\text{-value} = [0.0214]$.
    \item \textbf{Результат:} Данные [соответствуют/не соответствуют] заданному распределению.
\end{itemize}

\section{Заключение}
В данной работе была реализована программная проверка гипотез. Использование критерия Манна-Уитни позволило подтвердить устойчивость выводов $t$-критерия. При проверке критерия Пирсона было выполнено условие $np_k \ge 5$ путем объединения малочисленных интервалов, что гарантирует достоверность результата. Ошибка первого рода в данных экспериментах означала бы ложное отклонение гипотезы о равенстве параметров или виде распределения при их фактической верности.

\end{document}